% =====================================================================
%  CSD493: Project-1 (Monsoon 2026) -- DETAILED Project Proposal
%  Shiv Nadar Institution of Eminence, Dept. of CSE
%  Behavioral-Diff Detection of Malicious Browser Extension Updates
%
%  Revised version. Key changes vs. the first draft:
%    * Base paper set to an IEEE-published, top-tier work (Bui et al.,
%      IEEE S&P 2023); removed the unverifiable "Stock et al." citation.
%    * Multi-browser / cross-engine coverage added (Chromium family + Firefox).
%    * CRXcavator (sunset by Cisco/Duo in 2023) removed as a data source;
%      data plan re-based on synthetic corpus + validated real incidents.
%    * Manifest V3 target made explicit; Tranco replaces deprecated Alexa.
%    * Trace-determinism (record/replay) and a transparent weighted baseline added.
%    * Explicit scope/phasing and a publishable-novelty section.
%  Compile with pdfLaTeX.
% =====================================================================
\documentclass[11pt]{article}

\usepackage[a4paper,top=1in,bottom=1in,left=1in,right=1in]{geometry}
\usepackage{lmodern}
\usepackage[T1]{fontenc}
\usepackage{amsmath,amssymb}
\usepackage{booktabs}
\usepackage{enumitem}
\usepackage{xcolor}
\usepackage{titlesec}
\usepackage{fancyhdr}
\usepackage[colorlinks=true,linkcolor=black,urlcolor=blue,citecolor=blue]{hyperref}

\definecolor{snublue}{RGB}{20,60,120}

% --- Section styling ---
\titleformat{\section}{\large\bfseries\color{snublue}}{\thesection}{0.6em}{}
\titleformat{\subsection}{\normalsize\bfseries}{\thesubsection}{0.6em}{}
\titlespacing*{\section}{0pt}{1.4ex plus 1ex minus .2ex}{1ex}

% --- Running header/footer ---
\pagestyle{fancy}
\fancyhf{}
\lhead{\footnotesize CSD493: Project-1 Proposal}
\rhead{\footnotesize Shiv Nadar IoE}
\cfoot{\thepage}
\renewcommand{\headrulewidth}{0.4pt}

\setlength{\parskip}{0.6em}
\setlength{\parindent}{0pt}

\newcommand{\dem}[1]{\textbf{#1}}

\begin{document}

% ========================= Title block ===============================
\begin{center}
    {\large\bfseries Shiv Nadar Institution of Eminence}\\[2pt]
    {\normalsize Department of Computer Science and Engineering}\\[14pt]
    {\LARGE\bfseries\color{snublue} Behavioral-Diff Detection of Malicious\\[4pt]
     Browser Extension Updates}\\[8pt]
    {\itshape A sandboxed, cross-browser runtime-analysis framework for supply-chain
     threat detection in browser extensions}\\[14pt]
\end{center}

\begin{center}
\begin{tabular}{@{}ll@{}}
\textbf{Group Members:} & Bind Pratap Singh (2310110084), Suhani Deepak Agrawat (2310110717)\\
\textbf{Programme:}     & B.Tech CSE, Batch 2023--27\\
\textbf{Project Advisor:} & Dr. Sweta Mishra\\
\textbf{Area of Specialization:} & Cybersecurity\\
\textbf{Date:}          & May 23, 2026\\
\end{tabular}
\end{center}

\vspace{0.5em}
\hrule
\vspace{1em}

% ============================ Abstract ===============================
\begin{center}\textbf{Abstract}\end{center}
\begin{quote}
Browser extensions are a persistent supply-chain risk: a trusted, widely-installed extension can
become malicious after a single silent background update, reaching hundreds of thousands of users
within hours. Static review fails to catch obfuscated payloads, delayed activation, and behaviours
that only surface at runtime. This project proposes a sandboxed \emph{behavioral-diff} framework
that executes two consecutive versions of the same extension under identical, automated browsing
conditions and computes a weighted delta across their runtime traces. To remove environmental noise,
the two versions are driven against \emph{record/replayed} page content so the only meaningful
difference is the extension itself. The system actively triggers dormant payloads via synthetic DOM
interaction; records network traffic, extension-API calls, DOM access, and storage access; and scores
the behavioural change with both a transparent weighted baseline and a Random Forest classifier
trained on a synthetic weaponisation corpus and validated on real 2024--25 supply-chain incidents.
Crucially, the framework is \emph{cross-browser}: it analyses Chromium-family browsers
(Chrome, Edge, Brave) and Firefox, enabling detection of update-time divergence within and across
browser engines. The result is an explainable, per-update diff report that is directly actionable for
security teams. The system is designed to be implementable within the semester by two final-year CSE
students, while leaving a clear novelty surface suitable for a short research paper.
\end{quote}

\vspace{0.5em}

% =========================== Motivation ==============================
\section{Motivation}
The browser-extension ecosystem runs on a fundamental \emph{trust contract}: once a user installs an
extension, the browser grants it elevated access to every page they visit. This trust contract is
\textbf{not renewed on each update}. An extension that passes a careful initial review automatically
receives the same privileges after every subsequent silent update, with no re-consent from the user
and no re-review by the store. The same dynamic exists in \emph{every} major browser.

\begin{quote}
\textbf{December 2024 --- Cyberhaven Supply-Chain Attack.} A threat actor compromised the developer
account of Cyberhaven, a legitimate data-loss-prevention extension with 400,000 active users. A
malicious update was pushed to the Chrome Web Store and silently distributed to all users within
hours. The payload exfiltrated session cookies and authentication tokens from targeted financial
portals --- entirely bypassing static review, permission checks, and endpoint security tools. The
same campaign is now known to have compromised dozens of additional Chromium extensions.
\end{quote}

Static code analysis and permission-manifest inspection cannot catch this class of attack: malicious
logic can be obfuscated, condition-gated, or time-delayed. There is currently no open, systematic,
\emph{version-aware} tool that detects such behavioural divergence between consecutive extension
releases, and existing academic dynamic-analysis tools analyse a \emph{single} version of an
extension on a \emph{single} browser. This project targets exactly that gap.

% ======================== Problem Statement ==========================
\section{Problem Statement}
Given two consecutive releases of the same browser extension --- version $N$ (trusted baseline) and
version $N{+}1$ (candidate update), packaged as \texttt{.crx} (Chromium) or \texttt{.xpi}
(Firefox) --- the system must decide whether the behavioural change between them constitutes a
security threat.

The meaningful signal is \emph{change}: new behaviours that did not exist in version $N$, weighted by
their security sensitivity. Formally, the system computes a behavioural delta vector
\[
    \Delta \;=\; f(T_{N+1}) - f(T_{N}),
\]
where $T_N$ and $T_{N+1}$ are runtime traces of each version collected under \emph{identical,
replayed} browsing conditions, and $f(\cdot)$ is a feature-extraction function. A scoring function
then maps $\Delta$ to a risk score and outputs an interpretable diff report explaining which
behavioural changes drove the alert. Because the same extension is analysed on multiple browsers,
the framework also exposes \emph{cross-engine} divergence (behaviour that appears only under a
particular browser's API surface), which is itself a known evasion technique.

% ====================== Proposed Methodology =========================
\section{Proposed Methodology}
The system is a five-stage analysis pipeline, parameterised by target browser.

\begin{enumerate}[leftmargin=1.4em,itemsep=0.3em]
\item \dem{Scope \& Sandbox Construction.} Each version of an extension is unpacked and loaded into
an isolated Docker container running a headless browser instance --- Chromium (covering Chrome, Edge,
Brave, which share the Manifest V3 \texttt{.crx} format and the \texttt{chrome.*} API) and Firefox
(\texttt{.xpi}, \texttt{browser.*} WebExtensions API). All outbound traffic is intercepted by
\texttt{mitmproxy} (with its CA certificate installed inside the container to terminate TLS).
Containers are network-egress-restricted so live malware cannot reach real victims.

\item \dem{Deterministic Browsing \& Payload Fuzzing.} To make the diff meaningful, both versions are
driven against \emph{the same page content} using a record-and-replay proxy, eliminating ad/nonce/
timestamp noise that would otherwise produce false positives. Over this fixed content, the framework
performs synthetic DOM fuzzing --- clicking login buttons, populating \emph{dummy} credential forms ---
across Tranco-ranked top sites, financial login pages, and webmail portals, to activate
condition-gated logic. (Tranco replaces the deprecated Alexa list.)

\item \dem{Runtime Trace Collection.} The sandbox records a trace spanning four dimensions:
(i) outbound network requests (destinations, payload shape, exfiltration patterns),
(ii) DOM access events (reads of form fields, cookies-in-DOM, injected nodes),
(iii) storage/cookie access, and (iv) extension-API calls (\texttt{chrome.*}\,/\,\texttt{browser.*},
including MV3 service-worker activity such as \texttt{declarativeNetRequest}).

\item \dem{Feature Extraction \& Risk Weighting.} Raw traces are normalised and converted into the
numerical feature vector $\Delta$. Features are weighted by security sensitivity to manage false
positives --- e.g., domain-reputation delta, new-destination count, DOM-access intensity, permission
delta, and exfiltration-pattern matches.

\item \dem{Scoring: Baseline + Classifier.} $\Delta$ is scored two ways for robustness and
explainability: (a) a \emph{transparent weighted-rule baseline}, and (b) a \emph{Random Forest}
classifier (with XGBoost as a secondary model). Because confirmed malicious \emph{update pairs} are
scarce in the wild, the classifier is trained predominantly on a \emph{synthetic weaponisation
corpus} --- clean extensions into which known malicious behaviours (cookie exfiltration, credential
scraping, beaconing) are injected to create labelled $(N, N{+}1)$ pairs --- and is then
\emph{validated} on held-out real incidents from the 2024--25 campaigns. The agreement between the
interpretable baseline and the learned model is itself reported as a confidence signal.
\end{enumerate}

% ================= Multi-browser / cross-engine ======================
\section{Multi-Browser and Cross-Engine Coverage}
A core requirement of this project is coverage across browser families rather than a single browser.

\begin{itemize}[leftmargin=1.4em,itemsep=0.2em]
\item \dem{Chromium family (primary).} Chrome, Edge, Brave, and Opera share the Manifest~V3
\texttt{.crx} format and the \texttt{chrome.*} API, so a single instrumentation back-end covers all
of them with minimal additional work. This is also the family hit by the Cyberhaven campaign.
\item \dem{Firefox (second engine).} Firefox uses the \texttt{.xpi} package and the \texttt{browser.*}
WebExtensions API on the Gecko engine. Supporting it requires a separate loader and API-instrumentation
layer but yields genuine \emph{cross-engine} coverage and lets the framework detect behaviour that an
extension exposes only under one engine.
\item \dem{Out of scope.} Safari App Extensions (macOS-only, distinct model) are explicitly excluded.
\end{itemize}

A shared, browser-agnostic trace schema lets the same feature-extraction and scoring stages operate
over both engines, so the cross-engine capability is an architectural property rather than duplicated
code.

% =================== Novelty & Contribution ==========================
\section{Novelty and Academic Contribution}
The project is positioned to balance engineering with a defensible research contribution. The
publishable novelty surface is:

\begin{itemize}[leftmargin=1.4em,itemsep=0.25em]
\item \dem{Update-aware threat formulation.} We analyse the \emph{temporal} supply-chain threat ---
behavioural divergence between consecutive versions --- rather than scoring an extension in isolation,
which is what existing dynamic-analysis tools (e.g., Hulk) and IEEE detection work do.
\item \dem{Cross-engine behavioral diffing.} To our knowledge, no open framework computes a normalised
behavioural diff for the \emph{same} extension across both Chromium and Gecko engines, surfacing
engine-specific evasion.
\item \dem{Determinism via record/replay.} Diffing two versions over identical, replayed page content
isolates the extension as the only variable --- a methodological contribution that directly attacks
the dominant false-positive source in runtime extension analysis.
\item \dem{Explainable, dual-scored output.} A per-update diff report combining an interpretable
weighted baseline with a learned classifier produces actionable intelligence (which behaviour changed
and why) rather than an opaque label.
\end{itemize}

If the empirical results hold (high detection on real incidents at low FPR on legitimate updates),
this is a natural fit for a short paper / workshop submission in browser or supply-chain security.

% =================== Implementation & division ======================
\section{Implementation and Division of Work}
The project relies on an open-source Python stack and divides naturally into two workstreams.

\dem{Student A --- Bind Pratap Singh (Infrastructure).}
\begin{itemize}[leftmargin=1.4em,itemsep=0.1em]
\item Docker isolation, egress control, and Playwright browser automation (Chromium + Firefox).
\item \texttt{mitmproxy} traffic interception, the record/replay layer, and DOM/API event instrumentation.
\item Browser-agnostic trace-serialisation pipeline and extension-API call logging.
\end{itemize}

\dem{Student B --- Suhani Deepak Agrawat (Data \& ML).}
\begin{itemize}[leftmargin=1.4em,itemsep=0.1em]
\item Synthetic-weaponisation corpus generation, real-incident dataset curation, labelling, and
feature engineering/weighting.
\item Weighted-baseline scorer and Random Forest / XGBoost training and cross-validation.
\item FastAPI backend service and a Streamlit dashboard rendering the per-update diff report.
\end{itemize}

% ========================= Evaluation ================================
\section{Evaluation}
\begin{center}
\begin{tabular}{@{}ll@{}}
\toprule
\textbf{Metric} & \textbf{Target / Method}\\
\midrule
Detection Rate (TPR)      & Held-out confirmed malicious update pairs from 2024--25 campaigns \\
False Positive Rate (FPR) & Legitimate updates from top-100 extensions sampled over time \\
Cross-Browser Consistency & Agreement of the diff verdict across Chromium and Firefox \\
Evasion Resistance        & Obfuscated / delayed-activation variants of known payloads \\
Per-pair Latency          & End-to-end analysis time from extension input to report output \\
Report Quality            & Qualitative: are flagged diff items accurate and actionable? \\
\bottomrule
\end{tabular}
\end{center}

% ===================== Scope and phasing =============================
\section{Scope, Phasing, and Risk}
Because the full pipeline is substantial for one semester, deliverables are phased so that a working,
demonstrable system exists by the end of CSD493, with the remainder forming the follow-on project.

\begin{itemize}[leftmargin=1.4em,itemsep=0.2em]
\item \dem{Project-1 (this semester) --- core, must-have:} sandbox + deterministic browsing + the
four-dimension trace pipeline on the \emph{Chromium family}; synthetic corpus; the weighted baseline
plus a first Random Forest; an end-to-end diff report on a small set of real incidents.
\item \dem{Phase 2 / stretch:} the Firefox (cross-engine) back-end, the Streamlit dashboard,
evasion-resistance experiments, and ML refinement toward a paper-grade evaluation.
\item \dem{Key risks \& mitigations:} (i) \emph{Scarcity of real malicious update pairs} --- mitigated
by training on the synthetic corpus and reserving real incidents for validation; (ii) \emph{trace
non-determinism} --- mitigated by record/replay; (iii) \emph{safe malware handling} --- mitigated by
egress-restricted containers and dummy credentials only.
\end{itemize}

% ========================= References ================================
\section{References}
\begin{thebibliography}{99}
\setlength{\itemsep}{0.25em}

\bibitem{bui2023}
D. Bui, B. Tang, and K. G. Shin, ``Detection of Inconsistencies in Privacy Practices of Browser
Extensions (ExtPrivA),'' \emph{IEEE Symposium on Security and Privacy (S\&P)}, 2023.
\textbf{[Base paper.]}
\url{https://ieeexplore.ieee.org/document/10179338/}

\bibitem{hulk2014}
A. Kapravelos, C. Grier, N. Chachra, C. Kruegel, G. Vigna, and V. Paxson, ``Hulk: Eliciting Malicious
Behavior in Browser Extensions,'' \emph{USENIX Security}, 2014.
\url{https://www.usenix.org/conference/usenixsecurity14/technical-sessions/presentation/kapravelos}

\bibitem{writ2024}
G. Vasiliadis \emph{et al.}, ``WRIT: Web Request Integrity and Attestation Against Malicious Browser
Extensions,'' \emph{IEEE Transactions on Dependable and Secure Computing (TDSC)}, 2024.
\url{https://ieeexplore.ieee.org/document/10273582/}

\bibitem{vex2010}
S. Bandhakavi \emph{et al.}, ``VEX: Vetting Browser Extensions for Security Vulnerabilities,''
\emph{USENIX Security}, 2010.
\url{https://www.usenix.org/legacy/event/sec10/tech/full_papers/Bandhakavi.pdf}

\bibitem{didivet2024}
A. Aggarwal \emph{et al.}, ``Did I Vet You Before? Assessing the Chrome Web Store Vetting Process
through Browser Extension Similarity,'' arXiv:2406.00374, 2024.
\url{https://arxiv.org/abs/2406.00374}

\bibitem{karami2021}
S. Karami \emph{et al.}, ``Awakening the Web's Sleeper Agents: Misusing Service Workers for Privacy
Leakage,'' \emph{NDSS}, 2021.
\url{https://www.ndss-symposium.org/ndss-paper/awakening-the-webs-sleeper-agents-misusing-service-workers-for-privacy-leakage/}

\bibitem{fass2019}
A. Fass \emph{et al.}, ``HideNoSeek: Camouflaging Malicious JavaScript in Benign ASTs,''
\emph{ACM CCS}, 2019.
\url{https://dl.acm.org/doi/10.1145/3319535.3363196}

\bibitem{singh2025}
S. Singh \emph{et al.}, ``A Study on Malicious Browser Extensions in 2025,'' arXiv:2503.04292, 2025.
(Preprint --- motivation reference.)
\url{https://arxiv.org/abs/2503.04292}

\bibitem{javasith2025}
``JavaSith: A Client-Side Framework for Analyzing Potentially Malicious Extensions in Browsers,
VS Code, and NPM Packages,'' arXiv:2505.21263, 2025.
\url{https://arxiv.org/abs/2505.21263}

\bibitem{tranco2019}
V. Le Pochat \emph{et al.}, ``Tranco: A Research-Oriented Top Sites Ranking Hardened Against
Manipulation,'' \emph{NDSS}, 2019.
\url{https://tranco-list.eu/}

\bibitem{kaspersky2023}
Kaspersky, ``Malicious extensions in the Chrome Web Store,'' 2023.
\url{https://kaspersky.co.in/blog/dangerous-chrome-extensions-87-million/25861/}

\end{thebibliography}

\end{document}
